\chapter{Микроскопическое обрезание неравновесного носителя}
\label{ch:v10-nonequilibrium-bath-microscopic-cutoff}

Клеточная дисперсия делает пространство импульсов компактным, но сама
компактность ещё не гарантирует физического происхождения масштаба
обрезания. Для условного трёхмерного комплекса зона Бриллюэна задаёт
\begin{equation}
 k_{\rm BZ}\ell_{\rm cell}=\pi.
\end{equation}
Максимум символа лапласиана достигается в углу зоны и равен $12$. Поэтому
операторная норма частотного носителя даёт
\begin{equation}
 \frac{\omega_{\rm UV}\ell_{\rm cell}}{v_g}=2\sqrt3,
 \qquad
 \tau_{\rm UV}=\omega_{\rm UV}^{-1}
 =\frac{\ell_{\rm cell}}{2\sqrt3\,v_g}.
\end{equation}

Независимо, 43-уровневая спектральная мера предыдущей линии удовлетворяет
\begin{equation}
 \Lambda_{43}\ell_{\rm cell}=42.
\end{equation}
Обе конструкции имеют общий безразмерный родитель:
\begin{equation}
 \frac{\Lambda_{43}}{k_{\rm BZ}}=\frac{42}{\pi}.
\end{equation}
При дополнительной типизации $v_g=c$ отношение спектральной энергии
$\hbar c\Lambda_{43}$ к угловой энергии $\hbar\omega_{\rm UV}$ равно
$7\sqrt3$.

На координатах
\begin{equation}
 u=k_{\rm BZ}\ell_{\rm cell},\quad
 w=\frac{\omega_{\rm UV}\ell_{\rm cell}}{v_g},\quad
 s=\Lambda_{43}\ell_{\rm cell}
\end{equation}
положительный родитель
\begin{equation}
 \mathcal P_{\rm UV}
 =\frac12\bigl[(u-\pi)^2+(w-2\sqrt3)^2+(s-42)^2\bigr]
\end{equation}
имеет единственный минимум $(\pi,2\sqrt3,42)$ и единичный гессиан.

Однако для логарифмических переменных
$(k_{\rm BZ},\omega_{\rm UV},\Lambda_{43},\ell_{\rm cell},v_g)$
карта условий имеет ранг/ядро $3/2$. Якорь $v_g$ оставляет ранг/ядро
$4/1$; только независимое значение $\ell_{\rm cell}$ закрывает карту до
$5/0$. Следовательно, построен общий безразмерный cutoff-носитель, а не
абсолютное ультрафиолетовое число.

\begin{theorem}[Общий клеточно-спектральный cutoff]
Условная клеточная дисперсия и 43-уровневая спектральная мера допускают
строго положительный общий родитель, фиксирующий отношения
$\pi$, $2\sqrt3$, $42$ и $42/\pi$ без подстановки целевого физического
масштаба.
\end{theorem}

\begin{nogocriterion}[Компактная зона не выбирает метр]
Преобразование $\ell_{\rm cell}\mapsto a\ell_{\rm cell}$ и всех волновых
чисел в $a^{-1}$ раз сохраняет безразмерный родитель. Поэтому зона
Бриллюэна не выводит абсолютное обрезание, пока физическое ребро ячейки не
получено независимо.
\end{nogocriterion}

\begin{protocol}\begin{enumerate}
\item условный родитель обрезания: \textbf{$9/9$};
\item безразмерный клеточно-спектральный мост: \textbf{$3/3$};
\item происхождение микроскопического носителя: \textbf{$0/1$};
\item происхождение абсолютного cutoff: \textbf{$0/1$};
\item следующий гейт: \textbf{родитель времени корреляции ванны}.
\end{enumerate}\end{protocol}

Машинный сертификат сохранён в
\path{s2t_v10_cell_birth_four_volume_nonequilibrium_bath_microscopic_carrier_cutoff_parent_admission_gate_results.json}.